\documentclass[exercice,a4]{thedocument}%
\usepackage{pstricks}%
\usepackage{multido}%
\usepackage{tabularx}%
\usepackage{scratch3}%
\usepackage{enumitem}%
\usepackage[french]{babel}%
\startdatakeys
    \source{25GENMATAS1}
    \BO{2026}
    \cycle{4}
    \competencesDuSocle{CA,RA}
    \objectifsApprentissage{C4A1aa,C4A2ad,C4B2cb,C4D2aa,C4D3ac}
    \motsCles{Scratch,programme,triangle équilatéral,angle}
\enddatakeys
\begin{document}%
\providecommand{\degres}{^\circ}%
Dans cet exercice, aucune justification n'est attendue

Une élève utilise un logiciel de programmation pour réaliser des dessins à partir d'un triangle équilatéral. Elle crée le bloc \og triangle \fg ci-contre.% 

\medskip

\begin{minipage}{0.62\linewidth}
\begin{enumerate}
\item Sur la copie, recopier et compléter les lignes 3 et 5 du bloc \og triangle \fg{} afin qu'il dessine un triangle équilatéral.
\end{enumerate}
\end{minipage}\hfill
\begin{minipage}{0.30\linewidth}
\begin{scratch}[scale=0.8,num blocks, num start=1,baseline=2]]
\initmoreblocks{ %
{definir triangle}}
\blockpen{stylo en position d'écriture}
\blockrepeat{répéter \ovalnum{\phantom{aa}} fois}
{
\blockmove{avancer de \ovalvariable{côté} pas}
\blockmove{tourner \turnleft{} de \ovalnum{\phantom{aa}} degrés}
}

\end{scratch}
\end{minipage}

\medskip

Elle utilise maintenant le bloc \og triangle \fg{} pour l'intégrer dans différents programmes.

\medskip

\begin{enumerate}[start=2]
\item Associer chaque programme au dessin qu'il permet de réaliser.

\emph{On indiquera sur la copie, le numéro du dessin et la lettre du programme associé.}
\end{enumerate}

\begin{center}
\begin{tabularx}{\linewidth}{|*{2}{>{\centering \arraybackslash}X|}*{1}{>{\centering \arraybackslash}X}}\cline{1-2}
\textbf{Programme A}&\textbf{Programme B}&~\\
\begin{scratch}[scale=0.9]
\blockinit{quand \greenflag est cliqué}
\blockmove{s’orienter à \ovalnum{90} degrés}
\blockvariable{mettre \ovalvariable{côté} à \ovalnum{20}}
\blockrepeat{répéter \ovalnum{4} fois}
{
\blockmoreblocks{triangle}
\blockmove{tourner \turnleft de \ovalnum{90} degrés}
}
\end{scratch}
&
\begin{scratch}[scale=0.8]
\blockinit{Quand \greenflag est cliqué}
\blockmove{s’orienter à \ovalnum{90} degrés}
\blockvariable{mettre \ovalvariable{côté} à \ovalnum{20}}
\blockrepeat{répéter \ovalnum{4} fois}
{
\blockmoreblocks{triangle}
\blockmove{avancer de \ovalvariable{côté} pas}
}
\end{scratch}& \emph{On rappelle que l'instruction}

\begin{scratch}\blockmove{s’orienter à \ovalnum{90} degrés} \end{scratch}
permet de s'orienter vers la droite\\ \cline{1-2}
\end{tabularx}
\end{center}

\medskip
\begin{tabularx}{\linewidth}{|*{3}{>{\centering \arraybackslash}X|}}\hline
\psset{unit=.4cm,arrowsize=2pt 3}
\begin{pspicture}(-3,-2)(5,2)
\def\tri{\pspolygon(-1,0)(1,0)(0,1.732)}
\multido{\n=-3.0+2.5}{4}{\rput(\n,0){\tri}}
\end{pspicture} &
\psset{unit=.4cm,arrowsize=2pt 3}
\begin{pspicture}(-3,-2)(3,2)
\pspolygon(0,0)(2,0)(1,1.732)\pspolygon(0,0)(0,2)(-1.732,1)
\pspolygon(0,0)(-2,0)(-1,-1.732)\pspolygon(0,0)(0,-2)(1.732,-1)
\pspolygon(0,0)(2,0)(1,1.732)
\end{pspicture}
&
\psset{unit=0.4cm,arrowsize=2pt 3}
\begin{pspicture}(-3,-2)(3,2)
\def\tri{\pspolygon(-1,0)(1,0)(0,1.732)}
\multido{\n=-3+2}{4}{\rput(\n,0){\tri}}
\end{pspicture}\\ \hline
Dessin 1 &Dessin 2 &Dessin 3 \\ \hline
\end{tabularx}

\bigskip

\begin{enumerate}[start=3]
\item On s'intéresse maintenant au programme ci-dessous. En prenant 1~cm pour 10 pas, construire sur la copie le dessin obtenu lorsque le programme s'exécute.
\end{enumerate}

\begin{center}
\begin{scratch}[scale=0.8]\blockinit{Quand \greenflag est cliqué}
\blockmove{s’orienter à \ovalnum{90} degrés}
\blockvariable{mettre \ovalvariable{côté} à \ovalnum{20}}
\blockrepeat{répéter \ovalnum{4} fois}
{\blockmoreblocks{triangle}
\blockvariable{mettre \ovalvariable{côté} à \ovalvariable{côté} * \ovalnum{2} }}
\end{scratch}
\end{center}
\showthesource
\end{document}
